% ============================================
% DOC-03
% Gerado de molambudos.md — build_miolo.py
% ============================================
\section*{\textbf{DOC-03} --- Laudo de Necropsia}
\markboth{DOC-03 --- Laudo de Necropsia}{}

\noindent \textit{Instituto Médico Legal de Barbacena --- 14 de junho de 1979 --- Dr. Álvaro Torres}

\noindent \textbf{Protocolo:} IML-BA 1979.06.12-1260
\textbf{Paciente:} Joaquim Antônio Correia (Registro Hospitalar 1.260)
\textbf{Data do Óbito:} 12 de junho de 1979
\textbf{Médico Legista:} Dr. Álvaro Torres, CRM-MG 3.117

\vspace{0.5em}\noindent\rule{\textwidth}{0.5pt}\vspace{0.5em}

\noindent \textbf{Instituto Médico Legal de Barbacena --- Laudo de Necropsia N° 1.260 --- 14 de junho de 1979}

\noindent \textbf{Paciente:} Joaquim Antônio Correia
\textbf{Registro:} 1.260
\textbf{Idade estimada:} 72 anos (± 5 anos)
\textbf{Naturalidade:} Quixeramobim, Ceará
\textbf{Data do óbito:} 12 de junho de 1979
\textbf{Data da necropsia:} 14 de junho de 1979
\textbf{Médico legista:} Dr. Álvaro Torres --- CRM-MG 3.117
\textbf{Solicitante:} Direção do Hospital Colônia de Barbacena

\vspace{0.5em}\noindent\rule{\textwidth}{0.5pt}\vspace{0.5em}

\subsection*{1. EXAME EXTERNO}

\subsubsection*{1.1 Dados Antropométricos}

\noindent - \textbf{Peso:} 42 kg (Índice de massa corporal: 16,0 kg/m² --- desnutrição crônica grave)
- \textbf{Estatura:} 1,62 m
- \textbf{Estado de nutrição:} Caquexia acentuada. Panículo adiposo ausente. Massa muscular globalmente reduzida. As saliências ósseas --- especialmente clavículas, cristas ilíacas e patelas --- projetam-se sob a pele, com proeminência óssea acentuada. A pele está ressecada, com pregueamento cutâneo persistente (sinal de desidratação crônica). As unhas apresentam linhas de Beau (sulcos transversais que indicam episódios repetidos de desnutrição aguda ao longo da vida).

\subsubsection*{1.2 Achados Oculares}

\noindent \textbf{Icterícia unilateral assimétrica --- globo ocular direito.}

\noindent A esclera do olho direito apresenta coloração amarelo-ouro intensa, incomum, que se estende da córnea à conjuntiva bulbar em tom homogêneo. A coloração não é compatível com icterícia hepática clássica (que seria bilateral e difusa) nem com xantelasmas palpebrais. A esclera do olho esquerdo é branca, com vascularização normal.

\noindent \textit{Nota do legista:} À palpação, o globo ocular direito apresenta consistência ligeiramente mais firme que o esquerdo. Não há sinais de trauma recente. Não há prótese. A esclera amarela não desbota à compressão digital.

\noindent \textit{Hipóteses:}
- Sequela de trauma ocular não tratado
- Infecção neonatal não documentada
- \textit{(Anotação marginal, caligrafia diferente: ``Depois da exumação, esta hipótese caiu por terra. O olho brilhava no escuro. Não há explicação médica para o fenômeno. --- Lúcia M.'')}

\subsubsection*{1.3 Lesões por Eletroconvulsoterapia (ECT)}

\noindent O crânio apresenta \textbf{marcas bilaterais de queimadura elétrica} nas regiões temporais, imediatamente anteriores ao trago de ambas as orelhas. As marcas são ovais, com aproximadamente 3 cm ×2 cm, de coloração acastanhada, com bordas ligeiramente elevadas e tecido cicatricial espesso. A pele nessas regiões apresenta alopecia permanente --- os folículos capilares foram destruídos pela passagem repetida de corrente elétrica.

\noindent \textit{Análise:} As marcas são compatíveis com aplicação reiterada de eletrochoque sem gel condutor, utilizando eletrodos metálicos de contato seco. A ausência de gel condutor --- prática padrão no Colônia até a década de 1960 --- produzia queimaduras de espessura parcial que, ao cicatrizar, destruíam as glândulas sudoríparas e os folículos capilares da região.

\noindent \textbf{Fraturas compatíveis com ECT:}
- \textbf{Fratura por compressão da vértebra T5} (quinta vértebra torácica), com perda de aproximadamente 30\textbackslash{}\% da altura do corpo vertebral. Cicatrização antiga, sem calo ósseo exuberante. Típica de contração muscular violenta e descontrolada durante crise convulsiva tônico-clônica generalizada --- o paciente estava deitado de costas durante o procedimento, sem contenção adequada, e a flexão forçada da coluna durante a fase tônica da convulsão comprimiu a vértebra até fraturá-la.
- \textbf{Fratura consolidada do colo do fêmur esquerdo}, com encurtamento de aproximadamente 2 cm do membro inferior esquerdo (observado na deambulação, segundo registros de enfermagem). A fratura ocorreu durante uma sessão de eletrochoque em 1948 --- o paciente, na fase clônica da convulsão, foi projetado para fora da maca e caiu sobre o quadril.
- \textbf{Fratura consolidada do ramo mandibular direito}, com discreta assimetria facial. Durante a crise convulsiva, os músculos masseteres se contraíram com força suficiente para fraturar a mandíbula. A ausência de contenção bucal (protetores ou roletes de borracha) contribuiu para a lesão.

\subsubsection*{1.4 Marcas de Lobotomia Transorbital (Procedimento de Freeman)}

\noindent \textbf{Ambos os cantos internos dos olhos} apresentam cicatrizes puntiformes, com aproximadamente 2 mm de diâmetro, localizadas na junção entre a carúncula lacrimal e a prega semilunar. As cicatrizes são bilaterais e simétricas.

\noindent \textit{Análise:} As cicatrizes são compatíveis com a introdução de leucótomo transorbital --- instrumento semelhante a um picador de gelo, com aproximadamente 15 cm de comprimento e 6 mm de diâmetro, com ponta romba e haste graduada. O procedimento, descrito por Walter Freeman em 1946, consistia em:

\noindent 1. Submeter o paciente a eletrochoque para induzir inconsciência (``anestesia por convulsão'')
2. Posicionar o leucótomo sobre a conjuntiva ocular, no canto interno do olho
3. Inclinar o instrumento 20 graus em relação à linha média
4. Percutir o leucótomo com um martelo cirúrgico até que atravessasse a fina lâmina óssea da órbita (aproximadamente 2--3 mm de espessura)
5. Avançar o instrumento até 5--7 cm no interior do crânio
6. Movimentar o cabo do leucótomo em leque (lateromedial e mediolateral) para seccionar as fibras de associação do lobo pré-frontal
7. Repetir do lado contralateral

\noindent \textit{Achado relevante:} Na lâmina cribriforme (teto da órbita), identificam-se dois pequenos orifícios de aproximadamente 4 mm de diâmetro, com bordas perfeitamente lisas e cicatrizadas, posicionados simetricamente a 2,5 cm da linha média. A distância entre os orifícios e a angulação das margens ósseas confirmam que o instrumento foi introduzido na direção do lobo frontal medial. \textbf{Os orifícios estão pérvios} --- não houve cicatrização óssea completa.

\noindent \textit{Anotação marginal legista:} ``Apesar dos orifícios evidentes no osso, o exame macroscópico do cérebro não revela cicatrizes corticais compatíveis com o procedimento. A lobotomia pode não ter sido completada, ou o instrumento pode ter seguido um trajeto anatômico que não lesou o tecido cerebral de forma detectável. Nota: o prontuário não contém autorização para lobotomia. O procedimento foi realizado sem registro formal.''

\subsubsection*{1.5 Marcas de Contenção Física}

\noindent \textbf{Punhos e tornozelos:} Faixas de hiperpigmentação circumferencial em ambos os punhos (3 cm de largura) e ambos os tornozelos (2,5 cm de largura). A pele nessas áreas apresenta textura áspera, com perda de elasticidade e ausência de pelos. As faixas são contínuas e uniformes, características de \textbf{pressão prolongada por algemas de couro ou camisa-de-força}.

\noindent \textit{Observação:} As faixas nos punhos têm profundidade suficiente para ter causado compressão do nervo mediano bilateralmente (síndrome do túnel do carpo por contenção). A contratura dos dedos em semiflexão (mão em garra) sugere que o paciente permaneceu contido por períodos superiores a 30 dias consecutivos em múltiplas ocasiões.

\noindent \textbf{Tórax e abdome:} Escoriações lineares horizontais na altura dos mamilos, bilaterais, compatíveis com atrito repetido de tiras de couro ou lona (camisa-de-força). As escoriações apresentam diferentes graus de cicatrização, indicando uso recorrente ao longo de décadas.

\noindent \textbf{Quadril e coxas:} Áreas de alopecia e fibrose dérmica nas faces anteriores de ambas as coxas, na altura da prega inguinal. Compatíveis com imobilização prolongada em cama, com o paciente deitado de costas e as pernas contidas por tiras de couro ou lençóis molhados (que, ao secar, apertavam-se contra a pele, causando isquemia tecidual).

\subsubsection*{1.6 Cicatrizes de Automutilação}

\noindent Múltiplas cicatrizes lineares em ambos os antebraços, face anterior e posterior. Distribuição simétrica, temporalidade variada (desde cicatrizes recentes, com menos de 6 meses, até cicatrizes antigas, com décadas). O padrão é consistente com arranhões profundos autoinfligidos --- o paciente utilizava as próprias unhas para cavar sulcos na pele, abertos e repetidos, até atingir o tecido subcutâneo.

\noindent \textit{Diferenciação:} As cicatrizes têm bordas irregulares e largura variável, características de arranhões (unhas), não de cortes (lâmina). Não há cicatrizes sugestivas de tentativa de suicídio por incisão venosa. O padrão é de \textbf{automutilação compulsiva, não suicida}, provavelmente associada a episódios de agitação psicomotora.

\noindent \textbf{Observação adicional:} As cicatrizes mais antigas apresentam sobreposição de tatuagem por impregnação de terra (partículas de sujeira incorporadas à derme durante a cicatrização), sugerindo que as lesões foram abertas e expostas a solo contaminado. \textbf{Não há registro de profilaxia antitetânica no prontuário do paciente.}

\subsubsection*{1.7 Dentição e Cavidade Oral}

\noindent \textbf{Ausência de 12 dentes:} Incisivos superiores e inferiores (8), primeiro molar inferior direito, segundo pré-molar superior esquerdo, canino inferior esquerdo, terceiro molar inferior direito. Os dentes remanescentes apresentam:

\noindent - Desgaste acentuado da superfície oclusal (atrito por bruxismo, comum em pacientes submetidos a ECT sem protetor bucal)
- Cáries generalizadas com exposição pulpar em 4 dentes
- Fratura dentária não tratada em 2 dentes (compatível com contração mandibular durante ECT)

\noindent \textbf{Língua:} Duas grandes cicatrizes lineares nas bordas laterais longitudinais, aproximadamente 3 cm cada. São cicatrizes de mordedura repetida durante crises convulsivas. A ausência de protetor bucal durante o eletrochoque fez com que o paciente mordesse a própria língua em múltiplas ocasiões, com força suficiente para atravessar o tecido muscular. As cicatrizes são profundas e queloidianas --- o tecido de cicatrização formou nódulos palpáveis na borda da língua.

\subsubsection*{1.8 Amputação Traumática}

\noindent \textbf{Amputação parcial da falange distal do dedo mínimo da mão esquerda.} Cicatrização antiga, com coto rombo e ausência de falange ungueal. Não há registro de trauma que justifique a amputação nos documentos do prontuário.

\noindent \textit{Anotação marginal:} ``Cicatrização por segunda intenção, sem sutura. Possível mordedura ou avulsão traumática. O paciente também arrancava as próprias unhas com os dentes (registro de Abr/1937). A falange pode ter sido arrancada pelo próprio paciente durante episódio de automutilação.''

\subsubsection*{1.9 Queimaduras Térmicas}

\noindent Na face anterior do tórax e abdome, observam-se \textbf{manchas de hipopigmentação} dispersas, com bordas irregulares, de aproximadamente 2 a 8 cm de diâmetro. São cicatrizes de queimaduras de segundo grau superficial, já epitelizadas.

\noindent \textit{Análise:} O padrão térmico é compatível com contato com água quente ou superfícies aquecidas. Os registros de enfermagem do Colônia mencionam a prática de ``banhos terapêuticos'' --- pacientes eram imersos em água quente (40--50°C) ou fria (10--15°C) por períodos prolongados, como forma de ``sedação'' ou ``estimulação''. A ausência de padrão geométrico nas queimaduras sugere escaldadura por imersão, não por contato com objeto aquecido.

\subsection*{2. EXAME INTERNO}

\subsubsection*{2.1 Causa Mortis}

\noindent \textbf{Parada cardiorrespiratória} secundária a:
- Insuficiência cardíaca congestiva
- Desnutrição crônica grave (caquexia cardíaca)
- Broncoaspiração (alimentos aspirados para vias aéreas inferiores --- compatível com disfagia por sequelas neurológicas)

\subsubsection*{2.2 Aparelho Cardiovascular}

\noindent Coração com peso reduzido (220 g --- esperado para o biotipo: 280--340 g). Atrofia marrom do miocárdio (pigmentação por lipofuscina, típica de desnutrição prolongada). As fibras cardíacas apresentam vacuolização citoplasmática e perda de estriação transversal. Aorta com aterosclerose leve, incompatível com a idade estimada (a desnutrição crônica retardou o processo aterosclerótico).

\subsubsection*{2.3 Aparelho Digestivo}

\noindent \textbf{Estômago:} Vazio. Ausência absoluta de conteúdo alimentar. As paredes gástricas estão adelgaçadas (espessura média de 2 mm), com mucosa atrófica e pregueamento diminuído. O paciente não se alimentava há pelo menos 24 horas antes do óbito.

\noindent \textbf{Intestinos:\textit{ Sinais de má absorção crônica --- as vilosidades intestinais estão achatadas e encurtadas (atrofia vilositária subtotal). Presença de }Ascaris lumbricoides} (lombrigas) em alta carga parasitária (17 espécimes contados no intestino delgado, 4 no cólon ascendente). Os parasitas, cada um com 15--25 cm de comprimento, formavam um emaranhado que ocupava aproximadamente 40\textbackslash{}\% do lúmen intestinal.

\noindent \textit{Observação:} A alta carga parasitária, associada à desnutrição crônica, contribuiu para o estado de caquexia. O paciente competia com os próprios parasitas pelos poucos nutrientes que conseguia ingerir.

\noindent \textbf{Fígado:} Peso reduzido (980 g), com esteatose moderada e sinais de desnutrição proteico-calórica. Sem sinais de cirrose ou hepatopatia que justifiquem a icterícia unilateral observada no exame externo.

\subsubsection*{2.4 Sistema Nervoso Central}

\noindent \textbf{Encéfalo:} Peso 1.120 g (dentro dos parâmetros esperados para idade e biotipo). Os hemisférios cerebrais são simétricos. As meninges são transparentes, sem sinais de meningite ou hemorragia subaracnóidea.

\noindent \textit{Análise específica --- lobotomia:} \textbf{Ausência de cicatrizes corticais macroscópicas} compatíveis com lobotomia transorbital ou pré-frontal. As áreas orbitofrontais e pré-frontais mediais foram dissecadas e examinadas cuidadosamente. Não há gliose, cavitação ou alteração de parênquima que indique secção de fibras nervosas. A via de acesso dos leucótomos (através do teto das órbitas) foi identificada, mas o trajeto parece ter sido interrompido antes de atingir o parênquima cerebral, ou o instrumento pode ter seguido um plano anatômico que não lesou tecido cerebral detectável.

\noindent \textit{Nota do legista:} ``Considerando a ausência de dano cerebral, é possível que o procedimento de lobotomia tenha sido interrompido antes da lesão do parênquima, ou que o paciente tenha sido submetido apenas à perfuração óssea (como ``treinamento'' de residentes) sem a secção das fibras frontais. Não há registro no prontuário que esclareça.''

\noindent \textit{Análise específica --- eletrochoque:} ``Não há evidência macroscópica de que os procedimentos de eletrochoque repetido (estimados em centenas de sessões ao longo de 50 anos de internação) tenham causado dano cerebral detectável. Ausência de gliose, atrofia lobar ou alterações vasculares significativas.''

\noindent \textit{Anotação marginal em letra cursiva, aparentemente mais recente:\textit{ ``A ausência de dano cerebral não é um alívio. É o mais perturbador de todos os achados. Significa que Joaquim não era violento por lesão cerebral. A violência era }dele\textit{. Ou era de }outra coisa}. Os procedimentos não deixaram marca no cérebro --- mas deixaram no corpo inteiro. As marcas do corpo contam uma história. O cérebro calado é a parte mais gritante do laudo. --- L.M.''

\subsubsection*{2.5 Coluna Vertebral}

\noindent \textbf{Fratura por compressão em T5} (conforme descrito no item 1.3). A medula espinhal no nível T5--T6 apresenta leve achatamento ântero-posterior, sem compressão significativa. A fratura foi do tipo ``estável'' e não causou déficit neurológico. \textbf{Fraturas por compressão adicionais em T3 e T8}, menores, não identificadas no exame radiológico prévio.

\noindent \textit{Observação:} As três vértebras fraturadas (T3, T5, T8) formam uma tríade clássica de fraturas por ECT --- ocorrem durante a fase tônica da convulsão, quando a contração simultânea dos músculos flexores e extensores da coluna gera forças de compressão axial que excedem a resistência óssea.

\subsubsection*{2.6 Sistema Musculoesquelético --- Achados Adicionais}

\noindent - \textbf{Artropatia dos ombros:} Luxação recidivante glenoumeral bilateral (ombros deslocados e recolocados em múltiplas ocasiões). Associada à manobra de contenção conhecida como ``chave de braço'' --- técnica utilizada pelos enfermeiros do Colônia para imobilizar pacientes agitados, que consiste em hiperabduzir e rodar externamente o braço até que a articulação estale.
- \textbf{Cicatriz em couro cabeludo:} Cicatriz linear de 6 cm na região occipital, compatível com ferimento contuso contundente. Possível queda durante crise convulsiva ou agressão física por parte de funcionários.
- \textbf{Deformidade dos pés:} Hálux valgo bilateral severo, com sobreposição do segundo dedo sobre o hálux em ambos os pés. Deformidade causada pelo uso prolongado de calçados inadequados (tamanho único, sem numeração --- ``chinelos do Colônia'', fabricados pelos próprios pacientes na oficina de sapataria do hospital).

\subsection*{3. SÍNTESE DOS PROCEDIMENTOS IDENTIFICADOS NO CORPO}

\begin{moltablethree}
\hline
\textbf{Procedimento} & \textbf{Evidência Física} & \textbf{Período} \\
\hline
Ele\-tro\-con\-vul\-so\-te\-ra\-pia sem anestesia & Queimaduras temporais; fraturas vertebrais T3/T5/T8; fratura mandibular; fratura de fêmur; cicatrizes linguais; perda dentária & 1925--1979 (54 anos) \\
Lobotomia transorbital (Freeman) & Orifícios na lâmina cribriforme bilateral; cicatrizes conjuntivais & Década de 1950 (sem registro) \\
Camisa-de-força / contenção prolongada & Faixas de hiperpigmentação em punhos e tornozelos; contratura digital; alopecia por pressão & 1920--1979 (59 anos) \\
Banhos terapêuticos (escaldadura) & Cicatrizes de queimadura térmica em tórax e abdome & 1930--1960 \\
Insulinoterapia (coma insulínico) & \textit{(Ausência de marcador --- não deixa cicatrizes físicas)} & 1940--1950 \\
Contenção com lençóis molhados & Fibrose dérmica em coxas e quadril & 1930--1979 \\
\hline
\end{moltablethree}

\noindent \textit{Nota final (manuscrito, adicionado posteriormente):} Corpo apresenta 59 anos de marcadores de intervenção institucional: ECT sem anestesia (queimaduras temporais bilaterais, fraturas vertebrais T3/T5/T8), lobotomia transorbital (perfuração de lâmina cribriforme bilateral), contenção mecânica (contratura digital, fibrose dérmica em coxas), banhos terapêuticos (cicatrizes térmicas em tórax e abdome). \textbf{Causa mortis:} parada cardiorrespiratória. \textbf{Causa contribuinte:} caquexia por desnutrição crônica (peso estimado: 38 kg). \textbf{Comorbidade:} alta carga parasitária (Ascaris lumbricoides, 21 espécimes). Não foram encontrados marcadores de doença orgânica que justifiquem 62 anos de institucionalização. A lesão mais significativa do paciente não é anatômica.

\vspace{0.5em}\noindent\rule{\textwidth}{0.5pt}\vspace{0.5em}

\noindent Dr. Álvaro Torres

\noindent Médico Legista --- IML Barbacena

\noindent CRM-MG 3.117

\vspace{0.5em}\noindent\rule{\textwidth}{0.5pt}\vspace{0.5em}

\noindent \textit{Adendo (manuscrito, aparentemente adicionado posteriormente):} ``Este laudo foi redigido em 14 de junho de 1979, antes da exumação. Os achados descritos na seção seguinte --- Exumação e Achados Especiais --- foram adicionados após a abertura da cova e a recuperação do diário. O signatário declara que não há explicação médica para os fenômenos descritos no item 4 (Exumação), limitando-se a relatar o que foi apresentado pelo coveiro e pelo material recuperado. O diário permanece sob custódia do IML, anexado a este processo. --- A.T.''

\noindent \textit{Nota adicional (caligrafia de L.M., protocolo 2026):} ``Fragmentos ósseos de Strigiformes (coruja-buraqueira) foram recuperados do mesmo estrato da cova, junto ao material do paciente. A presença de restos de coruja no sepultamento de Joaquim foi registrada pelo coveiro como 'a ave que não saiu de perto'. Não foi encontrada explicação tafonômica para a deposição. Ossos de coruja não são achado comum em sepultamentos humanos no Colônia.''

\noindent\textit{\textbf{↪ Rotas:} \rota{MEM-04} \rotasep \rota{DOC-02} \rotasep \rota{LUC-03} \rotasep \rota{CONT-03}}